\documentclass[10pt,a4paper]{amsart}
\usepackage{amsmath,amssymb,amsthm,graphicx}
% Retained preamble for the Stacks Project chapter snapshot in this directory.
% The chapter's own preamble.tex is a copy of the shared Stacks Project preamble; the block
% below is sliced MECHANICALLY, line for line, from it. Nothing was added, removed or
% redefined; the wrapper class, the xr-hyper/hyperref wiring and the externaldocument list
% are omitted because the assembler supplies the document class.
%
% lines taken from preamble.tex: 3, 5, 9, 11, 12, 156, 157, 158, ... (51 lines total)
%
% package coverage restored from the chapter's own preamble (2026-10-13): xy, xy-2cell

\usepackage[all]{xy}
\xyoption{2cell}
\UseAllTwocells

\newenvironment{reference}{\comment}{\endcomment}
\newenvironment{slogan}{\comment}{\endcomment}
\newenvironment{history}{\comment}{\endcomment}
\theoremstyle{plain}
\newtheorem{theorem}[subsection]{Theorem}
\newtheorem{proposition}[subsection]{Proposition}
\newtheorem{lemma}[subsection]{Lemma}
\theoremstyle{definition}
\newtheorem{definition}[subsection]{Definition}
\newtheorem{example}[subsection]{Example}
\newtheorem{exercise}[subsection]{Exercise}
\newtheorem{situation}[subsection]{Situation}
\theoremstyle{remark}
\newtheorem{remark}[subsection]{Remark}
\newtheorem{remarks}[subsection]{Remarks}
\numberwithin{equation}{subsection}
\def\lim{\mathop{\mathrm{lim}}\nolimits}
\def\colim{\mathop{\mathrm{colim}}\nolimits}
\def\Spec{\mathop{\mathrm{Spec}}}
\def\Hom{\mathop{\mathrm{Hom}}\nolimits}
\def\Ext{\mathop{\mathrm{Ext}}\nolimits}
\def\SheafHom{\mathop{\mathcal{H}\!\mathit{om}}\nolimits}
\def\SheafExt{\mathop{\mathcal{E}\!\mathit{xt}}\nolimits}
\def\Sch{\mathit{Sch}}
\def\Mor{\mathop{\mathrm{Mor}}\nolimits}
\def\Ob{\mathop{\mathrm{Ob}}\nolimits}
\def\Sh{\mathop{\mathit{Sh}}\nolimits}
\def\NL{\mathop{N\!L}\nolimits}
\def\CH{\mathop{\mathrm{CH}}\nolimits}
\def\proetale{{pro\text{-}\acute{e}tale}}
\def\etale{{\acute{e}tale}}
\def\QCoh{\mathit{QCoh}}
\def\Ker{\mathop{\mathrm{Ker}}}
\def\Im{\mathop{\mathrm{Im}}}
\def\Coker{\mathop{\mathrm{Coker}}}
\def\Coim{\mathop{\mathrm{Coim}}}
\DeclareMathSymbol{\boxtimes}{\mathbin}{AMSa}{"02}
\def\QCohstack{\mathcal{QC}\!\mathit{oh}}
\def\Cohstack{\mathcal{C}\!\mathit{oh}}
\def\Spacesstack{\mathcal{S}\!\mathit{paces}}
\def\Quotfunctor{\mathrm{Quot}}
\def\Hilbfunctor{\mathrm{Hilb}}
\def\Curvesstack{\mathcal{C}\!\mathit{urves}}
\def\Polarizedstack{\mathcal{P}\!\mathit{olarized}}
\def\Complexesstack{\mathcal{C}\!\mathit{omplexes}}
\def\Pic{\mathop{\mathrm{Pic}}\nolimits}
\def\Picardstack{\mathcal{P}\!\mathit{ic}}
\def\Picardfunctor{\mathrm{Pic}}
\def\Deformationcategory{\mathcal{D}\!\mathit{ef}}
\usepackage{cleveref}
\crefname{originalproposition}{original proposition}{original propositions}
\Crefname{originalproposition}{Original proposition}{Original propositions}
\makeatletter
\newlabel{categories-example-presheaf}{{04TM}{source}{}{sourcepointer-1}{}}
\newlabel{categories-example-presheaf@cref}{{[originalproposition][0][]04TM}{[1][0][]source}{}{}{}}
\makeatother
\crefname{originalproposition}{original proposition}{original propositions}
\Crefname{originalproposition}{Original proposition}{Original propositions}
\makeatletter
\newlabel{formal-spaces-definition-affine-formal-spectrum}{{0AIF}{source}{}{sourcepointer-2}{}}
\newlabel{formal-spaces-definition-affine-formal-spectrum@cref}{{[originalproposition][0][]0AIF}{[1][0][]source}{}{}{}}
\makeatother
\crefname{originalproposition}{original proposition}{original propositions}
\Crefname{originalproposition}{Original proposition}{Original propositions}
\makeatletter
\newlabel{formal-spaces-lemma-factor-through-thickening}{{0AIA}{source}{}{sourcepointer-3}{}}
\newlabel{formal-spaces-lemma-factor-through-thickening@cref}{{[originalproposition][0][]0AIA}{[1][0][]source}{}{}{}}
\makeatother
\crefname{originalproposition}{original proposition}{original propositions}
\Crefname{originalproposition}{Original proposition}{Original propositions}
\makeatletter
\newlabel{more-algebra-definition-topological-ring}{{07E8}{source}{}{sourcepointer-4}{}}
\newlabel{more-algebra-definition-topological-ring@cref}{{[originalproposition][0][]07E8}{[1][0][]source}{}{}{}}
\makeatother
\crefname{originalproposition}{original proposition}{original propositions}
\Crefname{originalproposition}{Original proposition}{Original propositions}
\makeatletter
\newlabel{sdga-lemma-cartesian-cofinal-subcategory}{{0GZE}{source}{}{sourcepointer-5}{}}
\newlabel{sdga-lemma-cartesian-cofinal-subcategory@cref}{{[originalproposition][0][]0GZE}{[1][0][]source}{}{}{}}
\makeatother
\crefname{originalproposition}{original proposition}{original propositions}
\Crefname{originalproposition}{Original proposition}{Original propositions}
\makeatletter
\newlabel{sdga-proposition-derived-complete}{{0GZK}{source}{}{sourcepointer-6}{}}
\newlabel{sdga-proposition-derived-complete@cref}{{[originalproposition][0][]0GZK}{[1][0][]source}{}{}{}}
\makeatother
\crefname{originalproposition}{original proposition}{original propositions}
\Crefname{originalproposition}{Original proposition}{Original propositions}
\makeatletter
\newlabel{sdga-section-qc-inverse-system}{{0GZG}{source}{}{sourcepointer-7}{}}
\newlabel{sdga-section-qc-inverse-system@cref}{{[originalproposition][0][]0GZG}{[1][0][]source}{}{}{}}
\makeatother
\title{Stacks, Sheaves on Algebraic Stacks: proposition-QC-derived-complete}
\author{The Stacks Project Authors}
\date{stacks-project:stacks-sheaves.tex:local-snapshot-20261012; GNU FDL 1.2 or later}
\begin{document}
\maketitle
\noindent\textit{Editorial note.} Statement span is lines 5136-5149 of sources/stacks-algebraic-stacks-sheaves/pinned-original-source.tex; proof span is lines 5151-5174, the complete author proof (22 body lines).\par\medskip
\noindent\textit{Editorial note.} Navigation stubs: 7 cited-but-not-included labels are supplied as non-load-bearing pointers with their pinned tags where the snapshot tag index has them.\par\medskip
\noindent\textit{Editorial note.} extra\_packages is empty on purpose: the retained preamble already loads xy with options and a second declaration clashes in the assembler wrapper.\par\medskip
\noindent\textit{Editorial note.} Text is retained verbatim, including the author's own typos and placeholders.\par\medskip
\noindent\textit{Editorial note.} Queue index 26 of queue\_algebraic-stacks-sheaves.json (pb=1107, stmt 5136-5149, proof 5151-5174).\par\medskip
\noindent\textit{Allowed standard prerequisites (editorial).}\begin{itemize}
\item Explicit prerequisite boundary. The cited inputs are prerequisites; the argument in the retained span is not replaced by any cited result.
\end{itemize}\medskip
\ifdefined\hypertarget\hypertarget{sourcepointer-1}{}\fi
\noindent\textit{Source reference.} Original proposition 04TM: sources/stacks-algebraic-stacks-sheaves/pinned-original-source.tex, source lines 5174--5174. Cited by the retained statement or proof but not part of the collected spans; supplied only as a navigation pointer. The argument does not depend on its content.\par\medskip
\ifdefined\hypertarget\hypertarget{sourcepointer-2}{}\fi
\noindent\textit{Source reference.} Original proposition 0AIF: sources/stacks-algebraic-stacks-sheaves/pinned-original-source.tex, source lines 5174--5174. Cited by the retained statement or proof but not part of the collected spans; supplied only as a navigation pointer. The argument does not depend on its content.\par\medskip
\ifdefined\hypertarget\hypertarget{sourcepointer-3}{}\fi
\noindent\textit{Source reference.} Original proposition 0AIA: sources/stacks-algebraic-stacks-sheaves/pinned-original-source.tex, source lines 5174--5174. Cited by the retained statement or proof but not part of the collected spans; supplied only as a navigation pointer. The argument does not depend on its content.\par\medskip
\ifdefined\hypertarget\hypertarget{sourcepointer-4}{}\fi
\noindent\textit{Source reference.} Original proposition 07E8: sources/stacks-algebraic-stacks-sheaves/pinned-original-source.tex, source lines 5174--5174. Cited by the retained statement or proof but not part of the collected spans; supplied only as a navigation pointer. The argument does not depend on its content.\par\medskip
\ifdefined\hypertarget\hypertarget{sourcepointer-5}{}\fi
\noindent\textit{Source reference.} Original proposition 0GZE: sources/stacks-algebraic-stacks-sheaves/pinned-original-source.tex, source lines 5174--5174. Cited by the retained statement or proof but not part of the collected spans; supplied only as a navigation pointer. The argument does not depend on its content.\par\medskip
\ifdefined\hypertarget\hypertarget{sourcepointer-6}{}\fi
\noindent\textit{Source reference.} Original proposition 0GZK: sources/stacks-algebraic-stacks-sheaves/pinned-original-source.tex, source lines 5174--5174. Cited by the retained statement or proof but not part of the collected spans; supplied only as a navigation pointer. The argument does not depend on its content.\par\medskip
\ifdefined\hypertarget\hypertarget{sourcepointer-7}{}\fi
\noindent\textit{Source reference.} Original proposition 0GZG: sources/stacks-algebraic-stacks-sheaves/pinned-original-source.tex, source lines 5174--5174. Cited by the retained statement or proof but not part of the collected spans; supplied only as a navigation pointer. The argument does not depend on its content.\par\medskip
\begin{proposition}
\label{proposition-QC-derived-complete}
Let $S$ be a scheme. Let $X = \text{Spf}(A)$ where $A$ is an
an adic Noetherian topological $S$-algebra with ideal of definition $I$, see
More on Algebra, Definition \ref{more-algebra-definition-topological-ring}
and Formal Spaces, Definition
\ref{formal-spaces-definition-affine-formal-spectrum}.
Let $p : \mathcal{X} \to (\Sch/S)_{fppf}$
the be category fibred in sets associated to the functor $X$, see
Categories, Example \ref{categories-example-presheaf}.
Then $\mathit{QC}(\mathcal{X})$ is canonically equivalent to the
category $D_{comp}(A, I)$ of objects of $D(A)$ which are
derived complete with respect to $I$.
\end{proposition}

\begin{proof}
Recall that $X = \colim \Spec(A/I^n)$ as an fppf sheaf.
An object of $\mathcal{X}_{affine}$ is the same thing
as an affine scheme $U = \Spec(R)$ with a given morphism $f : U \to X$.
By Formal Spaces, Lemma \ref{formal-spaces-lemma-factor-through-thickening}
there exists an $n \geq 1$ such that
$f$ factors through the monomorphism $\Spec(A/I^n) \to X$.
Consider the full subcategory
$\mathcal{C} \subset \mathcal{X}_{affine}$ consisting of
the objects $\Spec(A/I^n) \to X$.
By the remarks just made and Differential Graded Sheaves, Lemma
\ref{sdga-lemma-cartesian-cofinal-subcategory}
restriction to $\mathcal{C}$ is an exact equivalence
$\mathit{QC}(\mathcal{X}) \to
\mathit{QC}(\mathcal{C}, \mathcal{O}|_\mathcal{C})$.
For simplicity, let us assume that
$I^n \not = I^{n + 1}$ for all $n \geq 1$.
Then $(\mathcal{C}, \mathcal{O}|_\mathcal{C})$ is isomorphic 
as a ringed site to the ringed site $(\mathbf{N}, (A/I^n))$, see
Differential Graded Sheaves, Section
\ref{sdga-section-qc-inverse-system}.
Hence we conclude by Differential Graded Sheaves, Proposition
\ref{sdga-proposition-derived-complete}.
\end{proof}

\end{document}
